\begin{textsample}{POS Dim 1 – 1998 – Score 21.00 – 1998/t000216.txt}  \label{ex:f1_pos_014}
David Gallo : This is Bill Lange.
I ' m Dave Gallo.
And we ' re going to tell you some stories from \textbf{the} \textbf{sea} here in video.
We ' ve got some of \textbf{the} most incredible video of Titanic that ' s ever been seen, and we ' re not going to show you any of it. \textbf{The} truth of \textbf{the} matter is that \textbf{the} Titanic — even though it ' s breaking all sorts of box office records — it ' s not \textbf{the} most exciting story from \textbf{the} \textbf{sea}.
And \textbf{the} problem, I think, is that we take \textbf{the} ocean for granted.
When you think about it, \textbf{the} oceans are 75 percent of \textbf{the} \textbf{planet}.
Most of \textbf{the} \textbf{planet} is ocean water. \textbf{The} average depth is about two miles.
Part of \textbf{the} problem, I think, is we stand at \textbf{the} beach, or we see images like this of \textbf{the} ocean, and you look out at this great big \textbf{blue} expanse, and it ' s shimmering and it ' s moving and there ' s waves and there ' s \textbf{surf} and there ' s tides, but you have no idea for what lies in there.
And in \textbf{the} oceans, there are \textbf{the} longest mountain ranges on \textbf{the} \textbf{planet}.
Most of \textbf{the} animals are in \textbf{the} oceans.
Most of \textbf{the} \textbf{earthquakes} and volcanoes are in \textbf{the} \textbf{sea}, at \textbf{the} bottom of \textbf{the} \textbf{sea}. \textbf{The} biodiversity and \textbf{the} biodensity in \textbf{the} ocean is higher, in places, than it is in \textbf{the} rainforests.
It ' s mostly unexplored, and yet there are beautiful sights like this that captivate us and make us become familiar with it.
But when you ' re standing at \textbf{the} beach, I want you to think that you ' re standing at \textbf{the} edge of a very unfamiliar world.
We have to have a very special technology to get into that unfamiliar world.
We use \textbf{the} submarine Alvin and we use cameras, and \textbf{the} cameras are something that Bill Lange has developed with \textbf{the} help of Sony.
Marcel Proust said,"\textbf{The} true voyage of discovery is not so much in seeking new landscapes as in having new eyes."People that have partnered with us have given us new eyes, not only on what exists — \textbf{the} new landscapes at \textbf{the} bottom of \textbf{the} \textbf{sea} — but also how we think about life on \textbf{the} \textbf{planet} itself.
Here ' s a jelly.
It ' s one of my favorites, because it ' s got all sorts of working parts.
This turns out to be \textbf{the} longest creature in \textbf{the} oceans.
It gets up to about 150 feet long.
But see all those different working things?
I love that kind of stuff.
It ' s got these \textbf{fishing} lures on \textbf{the} bottom.
They ' re going up and down.
It ' s got tentacles dangling, swirling around like that.
It ' s a colonial animal.
These are all individual animals banding together to make this one creature.
And it ' s got these jet thrusters up in front that it ' ll use in a moment, and a little light.
If you take all \textbf{the} big \textbf{fish} and schooling \textbf{fish} and all that, put them on one side of \textbf{the} scale, put all \textbf{the} jelly-type of animals on \textbf{the} other side, those guys win hands down.
Most of \textbf{the} biomass in \textbf{the} ocean is made out of creatures like this.
Here ' s \textbf{the} X-wing death jelly. \textbf{The} bioluminescence — they use \textbf{the} lights for attracting \textbf{mates} and attracting prey and communicating.
We couldn ' t begin to show you our archival stuff from \textbf{the} jellies.
They come in all different sizes and shapes.
Bill Lange : We tend to forget about \textbf{the} fact that \textbf{the} ocean is miles deep on average, and that we ' re real familiar with \textbf{the} animals that are in \textbf{the} first 200 or 300 feet, but we ' re not familiar with what exists from there all \textbf{the} way down to \textbf{the} bottom.
And these are \textbf{the} types of animals that live in that three- dimensional space, that micro-gravity environment that we really haven ' t explored.
You hear about giant squid and things like that, but some of these animals get up to be approximately 140, 160 feet long.
They ' re very little understood.
DG : This is one of them, another one of our favorites, because it ' s a little octopod.
You can actually see through his head.
And here he is, flapping with his ears and very gracefully going up.
We see those at all depths and even at \textbf{the} greatest depths.
They go from a couple of \textbf{inches} to a couple of feet.
They come right up to \textbf{the} submarine — they ' ll put their eyes right up to \textbf{the} window and peek inside \textbf{the} sub.
This is really a world within a world, and we ' re going to show you two.
In this case, we ' re passing down through \textbf{the} mid-ocean and we see creatures like this.
This is kind of like an undersea rooster.
This guy, that looks incredibly formal, in a way.
And then one of my favorites.
What a face!
This is basically scientific data that you ' re looking at.
It ' s footage that we ' ve collected for scientific purposes.
And that ' s one of \textbf{the} things that Bill ' s been doing, is providing scientists with this first view of animals like this, in \textbf{the} world where they belong.
They don ' t catch them in a net.
They ' re actually looking at them down in that world.
We ' re going to take a joystick, sit in front of our \textbf{computer}, on \textbf{the} Earth, and press \textbf{the} joystick forward, and fly around \textbf{the} \textbf{planet}.
We ' re going to look at \textbf{the} mid-ocean ridge, a 40, 000-mile long mountain range. \textbf{The} average depth at \textbf{the} top of it is about a mile and a half.
And we ' re over \textbf{the} Atlantic — that ' s \textbf{the} ridge right there — but we ' re going to go across \textbf{the} Caribbean, Central America, and end up against \textbf{the} Pacific, nine degrees north.
We make maps of these mountain ranges with sound, with sonar, and this is one of those mountain ranges.
We ' re coming around a cliff here on \textbf{the} right. \textbf{The} height of these mountains on either side of this valley is greater than \textbf{the} Alps in most cases.
And there ' s tens of thousands of those mountains out there that haven ' t been mapped yet.
This is a volcanic ridge.
We ' re getting down further and further in scale.
And eventually, we can come up with something like this.
This is an icon of our robot, Jason, it ' s called.
And you can sit in a room like this, with a joystick and a headset, and drive a robot like that around \textbf{the} bottom of \textbf{the} ocean in real time.
One of \textbf{the} things we ' re trying to do at Woods Hole with our partners is to bring this virtual world — this world, this unexplored region — back to \textbf{the} laboratory.
Because we see it in bits and pieces right now.
We see it either as sound, or we see it as video, or we see it as photographs, or we see it as chemical sensors, but we never have yet put it all together into one interesting picture.
Here ' s where Bill ' s cameras really do shine.
This is what ' s called a hydrothermal vent.
And what you ' re seeing here is a cloud of densely packed, hydrogen-sulfide-rich water coming out of a volcanic axis on \textbf{the} \textbf{sea} floor.
Gets up to 600, 700 degrees F, somewhere in that range.
So that ' s all water under \textbf{the} \textbf{sea} — a mile and a half, two miles, three miles down.
And we knew it was volcanic back in \textbf{the} ' 60s, ' 70s.
And then we had some hint that these things existed all along \textbf{the} axis of it, because if you ' ve got volcanism, water ' s going to get down from \textbf{the} \textbf{sea} into cracks in \textbf{the} \textbf{sea} floor, come in contact with magma, and come shooting out hot.
We weren ' t really aware that it would be so rich with sulfides, hydrogen sulfides.
We didn ' t have any idea about these things, which we call chimneys.
This is one of these hydrothermal vents.
Six hundred degree F water coming out of \textbf{the} Earth.
On either side of us are mountain ranges that are higher than \textbf{the} Alps, so \textbf{the} setting here is very dramatic.
BL : \textbf{The} white material is a type of bacteria that thrives at 180 degrees C.
DG : I think that ' s one of \textbf{the} greatest stories right now that we ' re seeing from \textbf{the} bottom of \textbf{the} \textbf{sea}, is that \textbf{the} first thing we see coming out of \textbf{the} \textbf{sea} floor after a volcanic eruption is bacteria.
And we started to wonder for a long time, how did it all get down there?
What we find out now is that it ' s probably coming from inside \textbf{the} Earth.
Not only is it coming out of \textbf{the} Earth — so, biogenesis made from volcanic activity — but that bacteria supports these colonies of life. \textbf{The} pressure here is 4, 000 \textbf{pounds} per square \textbf{inch}.
A mile and a half from \textbf{the} \textbf{surface} to two miles to three miles — no sun has ever gotten down here.
All \textbf{the} energy to support these life forms is coming from inside \textbf{the} Earth — so, chemosynthesis.
And you can see how dense \textbf{the} population is.
These are called tube worms.
BL : These worms have no digestive system.
They have no mouth.
But they have two types of gill structures.
One for extracting oxygen out of \textbf{the} deep-sea water, another one which houses this chemosynthetic bacteria, which takes \textbf{the} hydrothermal fluid — that hot water that you saw coming out of \textbf{the} bottom — and converts that into simple sugars that \textbf{the} tube worm can digest.
DG : You can see, here ' s a \textbf{crab} that lives down there.
He ' s managed to grab a tip of these worms.
Now, they normally retract as soon as a \textbf{crab} touches them.
Oh!
Good going.
So, as soon as a \textbf{crab} touches them, they retract down into their shells, just like your fingernails.
There ' s a whole story being played out here that we ' re just now beginning to have some idea of because of this new camera technology.
BL : These worms live in a real temperature extreme.
Their foot is at about 200 degrees C and their head is out at three degrees C, so it ' s like having your hand in boiling water and your foot in freezing water.
That ' s how they like to live.
DG : This is a female of this kind of worm.
And here ' s a male.
You watch.
It doesn ' t take long before two guys here — this one and one that will show up over here — start to fight.
Everything you see is played out in \textbf{the} pitch black of \textbf{the} deep \textbf{sea}.
There are never any lights there, except \textbf{the} lights that we bring.
Here they go.
On one of \textbf{the} last dive series, we counted 200 species in these areas — 198 were new, new species.
BL : One of \textbf{the} big problems is that for \textbf{the} biologists working at these sites, it ' s rather difficult to collect these animals.
And they disintegrate on \textbf{the} way up, so \textbf{the} imagery is critical for \textbf{the} science.
DG : Two octopods at about two miles depth.
This pressure thing really amazes me — that these animals can exist there at a depth with pressure enough to crush \textbf{the} Titanic like an empty Pepsi can.
What we saw up till now was from \textbf{the} Pacific.
This is from \textbf{the} Atlantic.
Even greater depth.
You can see this \textbf{shrimp} is harassing this poor little guy here, and he ' ll bat it away with his claw.
Whack!
And \textbf{the} same thing ' s going on over here.
What they ' re getting at is that — on \textbf{the} back of this \textbf{crab} — \textbf{the} foodstuff here is this very strange bacteria that lives on \textbf{the} backs of all these animals.
And what these \textbf{shrimp} are trying to do is actually harvest \textbf{the} bacteria from \textbf{the} backs of these animals.
And \textbf{the} \textbf{crabs} don ' t like it at all.
These long filaments that you see on \textbf{the} back of \textbf{the} \textbf{crab} are actually created by \textbf{the} product of that bacteria.
So, \textbf{the} bacteria grows hair on \textbf{the} \textbf{crab}.
On \textbf{the} back, you see this again. \textbf{The} red dot is \textbf{the} laser light of \textbf{the} submarine Alvin to give us an idea about how far away we are from \textbf{the} vents.
Those are all \textbf{shrimp}.
You see \textbf{the} hot water over here, here and here, coming out.
They ' re clinging to a \textbf{rock} face and actually scraping bacteria off that \textbf{rock} face.
Here ' s a tiny, little vent that ' s come out of \textbf{the} side of that pillar.
Those pillars get up to several stories.
So here, you ' ve got this valley with this incredible alien landscape of pillars and hot \textbf{springs} and volcanic eruptions and \textbf{earthquakes}, inhabited by these very strange animals that live only on chemical energy coming out of \textbf{the} ground.
They don ' t need \textbf{the} sun at all.
BL : You see this white V-shaped mark on \textbf{the} back of \textbf{the} \textbf{shrimp}?
It ' s actually a light-sensing organ.
It ' s how they find \textbf{the} hydrothermal vents. \textbf{The} vents are emitting a black body radiation — an IR signature — and so they ' re able to find these vents at considerable distances.
DG : All this stuff is happening along that 40, 000-mile long mountain range that we ' re calling \textbf{the} ribbon of life, because just even today, as we speak, there ' s life being generated there from volcanic activity.
This is \textbf{the} first time we ' ve ever tried this any place.
We ' re going to try to show you high definition from \textbf{the} Pacific.
We ' re moving up one of these pillars.
This one ' s several stories tall.
In it, you ' ll see that it ' s a \textbf{habitat} for a lot of different animals.
There ' s a funny kind of hot \textbf{plate} here, with vent water coming out of it.
So all of these are individual homes for worms.
Now here ' s a closer view of that community.
Here ' s \textbf{crabs} here, worms here.
There are smaller animals crawling around.
Here ' s pagoda structures.
I think this is \textbf{the} neatest-looking thing.
I just can ' t get over this — that you ' ve got these little chimneys sitting here smoking away.
This stuff is toxic as hell, by \textbf{the} way.
You could never get a permit to dump this in \textbf{the} ocean, and it ' s coming out all from it.
It ' s unbelievable.
It ' s basically sulfuric acid, and it ' s being just dumped out, at incredible rates.
And animals are thriving — and we probably came from here.
That ' s probably where we \textbf{evolved} from.
BL : This bacteria that we ' ve been talking about turns out to be \textbf{the} most simplest form of life found.
There are a number of groups that are proposing that life \textbf{evolved} at these vent sites.
Although \textbf{the} vent sites are short-lived — an individual site may last only 10 years or so — as an ecosystem they ' ve been stable for millions — well, billions — of years.
DG : It works too well.
You see there ' re some \textbf{fish} inside here as well.
There ' s a \textbf{fish} sitting here.
Here ' s a \textbf{crab} with his claw right at \textbf{the} end of that tube worm, waiting for that worm to stick his head out.
BL : \textbf{The} biologists right now cannot explain why these animals are so active. \textbf{The} worms are growing \textbf{inches} per week!
DG : I already said that this site, from a human perspective, is toxic as hell.
Not only that, but on top — \textbf{the} lifeblood — that plumbing system turns off every year or so.
Their plumbing system turns off, so \textbf{the} sites have to move.
And then there ' s \textbf{earthquakes}, and then volcanic eruptions, on \textbf{the} order of one every five years, that completely \textbf{wipes} \textbf{the} area out.
Despite that, these animals grow back in about a year ' s time.
You ' re talking about biodensities and biodiversity, again, higher than \textbf{the} rainforest that just \textbf{springs} back to life.
Is it sensitive?
Yes.
Is it fragile?
No, it ' s not really very fragile.
I ' ll end up with saying one thing.
There ' s a story in \textbf{the} \textbf{sea}, in \textbf{the} waters of \textbf{the} \textbf{sea}, in \textbf{the} sediments and \textbf{the} \textbf{rocks} of \textbf{the} \textbf{sea} floor.
It ' s an incredible story.
What we see when we look back in time, in those sediments and \textbf{rocks}, is a record of Earth history.
Everything on this \textbf{planet} — everything — works by cycles and rhythms. \textbf{The} \textbf{continents} move apart.
They come back together.
Oceans come and go.
Mountains come and go.
Glaciers come and go.
El Nino comes and goes.
It ' s not a disaster, it ' s rhythmic.
What we ' re learning now, it ' s almost like a symphony.
It ' s just like music — it really is just like music.
And what we ' re \textbf{learning} now is that you can ' t listen to a five-billion-year long symphony, get to today and say,"Stop!
We want tomorrow ' s note to be \textbf{the} same as it was today."It ' s absurd.
It ' s just absurd.
So, what we ' ve got to learn now is to find out where this \textbf{planet} ' s going at all these different scales and work with it.
Learn to manage it. \textbf{The} concept of preservation is futile.
Conservation ' s tougher, but we can probably get there.
Thank you very much.
Thank you.

% matched lemmas: blue, computer, continent, crab, earthquake, evolve, fish, fishing, habitat, inch, learning, mate, planet, plate, pound, rock, sea, shrimp, spring, surf, surface, the, wipe
\end{textsample}
